%(BEGIN_QUESTION)
% Copyright 2009, Tony R. Kuphaldt, released under the Creative Commons Attribution License (v 1.0)
% This means you may do almost anything with this work of mine, so long as you give me proper credit

Du skal forberede deg til klasseromsøkten (teori) ved å gjennomføre leseoppgaver og/eller forsøke å besvare oppgavespørsmål for hver dag før du møter i timen. Dette krever betydelig selvstendig research og problemløsning.

Hva bør du gjøre hvis du møter et spørsmål som fullstendig forvirrer deg, og du ikke aner hvordan det skal besvares? Og tilsvarende: hva bør du gjøre hvis du møter en del av pensumteksten du rett og slett ikke klarer å forstå?

\underbar{file i00122}
%(END_QUESTION)





%(BEGIN_ANSWER)

Hvis du står helt fast i et spørsmål eller i en del av pensumteksten, selv etter at du har brukt tilgjengelig studietid før timen, bør du markere disse punktene tydelig i notatene dine og be om hjelp med en gang ved starten av timen. Sjansen er stor for at du ikke er den eneste med akkurat det spørsmålet, og at spørsmålet ditt også hjelper andre.

%(END_ANSWER)





%(BEGIN_NOTES)


%INDEX% Course organization, assessment: preparation for theory sessions

%(END_NOTES)
